\section{When Step 1 fails: long-range scattering and infrared sectors}
\label{sec:long-range-ir}
% BEGIN CURATED SUBJECT INDEX
\index{wave operator}
\index{scattering!Coulomb}
\index{Dollard comparison dynamics}
\index{scattering!three-body Coulomb}
\index{infrared problem}
\index{dressing state}
% END CURATED SUBJECT INDEX

The standard scattering construction compares the full dynamics with a free
reference at large times.  This works for a short-range interaction because
the phase accumulated on an escaping classical trajectory is finite.  If
$\bm r(t)\simeq\bm v t$ and $V(r)=O(r^{-1-\epsilon})$, then
\begin{equation}
  \int^{\pm\infty}\lvert V(\bm v t)\rvert\dd t<\infty .
  \label{eq:short-range-trajectory-integral}
\end{equation}
For a Coulomb tail, $V_C(r)=g/r$, the same integral grows as
$g\lvert\bm v\rvert^{-1}\ln|t|$.  The interaction never becomes negligible
in phase, even though the force tends to zero.  Consequently the ordinary
wave operators comparing $H$ directly with $H_0$ do not exist in the form
used for short-range scattering.  The cure is not to discard the phase but to
put it into the asymptotic dynamics \cite{Dollard:1964,Dollard:1971,
DerezinskiGerard:1997}.

This observation links three problems that are often taught separately.  The
two-body Coulomb problem requires Coulomb-distorted plane waves.  The
three-body Coulomb problem requires different correlated expansions in
different asymptotic sectors.  QED requires coherent soft-photon dressings and,
in general, representations inequivalent to the vacuum Fock representation.
In all three cases a long-range correlation changes the definition of an
asymptotic state.

\subsection{Two-body Coulomb waves rather than free plane waves}
\label{subsec:two-body-coulomb}
% BEGIN CURATED SUBJECT INDEX
\index{Jost function}
\index{Lippmann--Schwinger equation}
\index{cross section}
\index{partial-wave expansion}
\index{boundary condition!incoming and outgoing}
\index{scattering!Coulomb}
% END CURATED SUBJECT INDEX

Consider
\begin{equation}
  H=H_0+\frac{g}{r},
  \qquad
  H_0=-\frac{\hbar^2}{2\mu}\nabla^2,
  \qquad
  E=\frac{\hbar^2k^2}{2\mu},
  \label{eq:coulomb-hamiltonian-scattering}
\end{equation}
and introduce the Sommerfeld parameter
\begin{equation}
  \eta=\frac{\mu g}{\hbar^2 k}
  =\frac{g}{\hbar v}.
  \label{eq:sommerfeld-parameter}
\end{equation}
For repulsion, $g>0$ and $\eta>0$.  A conventionally normalized outgoing
solution is
\begin{equation}
  \psi_{\bm k}^{C(+)}(\bm r)
  =\rme^{-\pi\eta/2}\Gamma(1+\rmi\eta)
  \rme^{\rmi\bm k\cdot\bm r}
  {}_1F_1\left(-\rmi\eta;1;
  \rmi(kr-\bm k\cdot\bm r)\right).
  \label{eq:coulomb-hypergeometric-wave}
\end{equation}
These formulas and their partial-wave counterparts are standard consequences
of the confluent hypergeometric representation of the Coulomb problem
\cite{Taylor:1972,PerelomovZeldovich:1998}.
At fixed nonforward angle and large $r$, its two pieces have the form
\begin{align}
  \psi_{\bm k}^{C(+)}(\bm r)
  &\sim
  \underbrace{\rme^{\rmi\left[
  \bm k\cdot\bm r
  \eta\ln(kr-\bm k\cdot\bm r)\right]}}
  _{\text{Coulomb-distorted incident wave}}
  \notag\\
  &\quad+
  f_C(\vartheta)
  \frac{\rme^{\rmi\left[
  kr-\eta\ln(2kr)\right]}}{r}.
  \label{eq:coulomb-full-asymptotic}
\end{align}
Thus the incident object is not a plane wave.  Both the incident and outgoing
parts contain logarithmic phases.  The expansion is nonuniform in the forward
direction, where the distinction between the two terms itself becomes
delicate.

With
\begin{equation}
  \sigma_l(\eta)=\arg\Gamma(l+1+\rmi\eta),
  \label{eq:coulomb-phase-shift}
\end{equation}
the Coulomb amplitude is
\begin{equation}
  f_C(\vartheta)
  =-\frac{\eta}{2k\sin^2(\vartheta/2)}
  \rme^{
  2\rmi\sigma_0
  -2\rmi\eta\ln\sin(\vartheta/2)},
  \label{eq:coulomb-amplitude}
\end{equation}
and its modulus gives the Rutherford cross section
\begin{equation}
  \frac{\rmd\sigma_C}{\rmd\Omega}
  =\frac{\eta^2}{4k^2\sin^4(\vartheta/2)} .
  \label{eq:rutherford-cross-section}
\end{equation}
The forward divergence is physical for an ideal unscreened infinite-range
force; a realistic angular resolution or screening prescription is needed for
an integrated rate.

The radial Coulomb functions make the modified boundary condition especially
clear.  With $\rho=kr$,
\begin{equation}
  H_l^{(\pm)}(\eta,\rho)
  \sim
  \rme^{
  \pm\rmi\left[
  \rho-\eta\ln(2\rho)-\frac{l\pi}{2}+\sigma_l
  \right]}.
  \label{eq:coulomb-H-asymptotic}
\end{equation}
If a short-range interaction $W$ is added to $g/r$, the regular solution can
be written asymptotically as
\begin{equation}
  u_l(r)\sim\frac{1}{2\rmi}
  \left[
  S_l^{N}(k)H_l^{(+)}(\eta,kr)
  -H_l^{(-)}(\eta,kr)
  \right],
  \label{eq:coulomb-modified-radial-S}
\end{equation}
where $S_l^{N}=\rme^{2\rmi\delta_l^N}$ is the Coulomb-modified nuclear
matrix.  Relative to free partial waves, the full matrix is
\begin{equation}
  S_l(k)=\rme^{2\rmi\sigma_l(\eta)}S_l^{N}(k).
  \label{eq:full-coulomb-nuclear-S}
\end{equation}
Likewise, the correct distorted-wave Lippmann--Schwinger equation uses the
Coulomb Hamiltonian $H_C=H_0+g/r$ as reference:
\begin{equation}
  \ket{\psi_{\bm k}^{(+)}}
  =\ket{\phi_{\bm k}^{C(+)}}
  +(E+\rmi0-H_C)^{-1}W\ket{\psi_{\bm k}^{(+)}} .
  \label{eq:coulomb-distorted-LS}
\end{equation}
Coulomb-modified Jost functions are coefficients of
$H_l^{(\pm)}$, not of $\rme^{\pm\rmi kr}$.  Factoring out the known
$\Gamma$-function and logarithmic structure is essential before searching for
additional short-range resonance poles.

\subsection{Dollard comparison dynamics}
\label{subsec:dollard}
% BEGIN CURATED SUBJECT INDEX
\index{cross section}
\index{boundary condition!incoming and outgoing}
\index{scattering!Coulomb}
\index{Dollard comparison dynamics}
% END CURATED SUBJECT INDEX

The same modification has a time-dependent formulation.  Along a free
trajectory the asymptotic Coulomb operator in momentum space is
\begin{equation}
  V_D(\bm p)=\frac{g\mu}{|\bm p|}.
  \label{eq:dollard-asymptotic-potential}
\end{equation}
Choose an arbitrary reference time $t_0>0$ and define
\begin{align}
  \Phi_D(t,\bm p)
  &=\frac{\operatorname{sgn}(t)}{\hbar}
  V_D(\bm p)\ln\frac{|t|}{t_0},
  \label{eq:dollard-phase}\\
  U_D(t)
  &=\rme^{-\rmi H_0t/\hbar}
  \rme^{-\rmi\Phi_D(t,\bm p)} .
  \label{eq:dollard-comparison-dynamics}
\end{align}
Up to sign conventions for incoming and outgoing labels, the modified wave
operators are
\begin{equation}
  \Omega_D^{(\pm)}
  =\operatorname*{s-lim}_{t\to\mp\infty}
  \rme^{\rmi Ht/\hbar}U_D(t)\vsub{P}{ac}(H_0).
  \label{eq:dollard-wave-operators}
\end{equation}
The logarithm removed by $U_D(t)$ is the time-domain version of the logarithm
in Eq.~\eqref{eq:coulomb-full-asymptotic}.  Changing $t_0$ changes a
conventional channel phase but not a cross section.  A screened Coulomb
potential can also be used, but its screening radius cannot simply be sent to
infinity at the level of the undressed $S$ matrix: a divergent phase must be
removed first.

For many particles, asymptotic velocities and cluster decompositions enter
the comparison dynamics.  Pairwise Coulomb phases are correlated because the
relative momenta share the same particle momenta.  This is the first sign that
the many-body problem is not obtained by multiplying independent copies of
the two-body answer.

\subsection{Three charged particles: known sectors and the missing uniform form}
\label{subsec:three-body-coulomb}
% BEGIN CURATED SUBJECT INDEX
\index{compact operator}
\index{scattering!Coulomb}
\index{scattering!three-body Coulomb}
\index{Jacobi coordinate}
\index{Faddeev equation}
\index{nuclear breakup}
% END CURATED SUBJECT INDEX

For three particles,
\begin{equation}
  H=H_0+\sum_{1\leq i<j\leq3}\frac{g_{ij}}{r_{ij}} .
  \label{eq:three-body-coulomb-H}
\end{equation}
Choose Jacobi coordinates $(\bm x_\alpha,\bm y_\alpha)$, where
$\bm x_\alpha$ is the separation inside the pair
$\alpha=(\beta\gamma)$ and $\bm y_\alpha$ separates the spectator from that
pair.  Infinity has several geometrically different sectors:
\begin{align}
  \Omega_0&:\quad
  r_{12}\sim r_{23}\sim r_{31}\to\infty,
  \label{eq:three-body-breakup-sector}\\
  \Omega_\alpha&:\quad
  y_\alpha\to\infty,
  \qquad x_\alpha/y_\alpha\to0 .
  \label{eq:three-body-cluster-sector}
\end{align}
The first is the democratic three-body breakup region.  The second is a
cluster region in which a pair remains close compared with the spectator
distance.  Their asymptotic expansions are not uniform limits of one another.

In $\Omega_0$, the Redmond or 3C incident wave incorporates the leading
Coulomb phase of every pair.  In a common notation,
\begin{align}
  \vsub{\Psi}{3C}(\bm X)
  &=
  \rme^{\rmi\bm K\cdot\bm X}
  \prod_{i<j}
  \left[
  \rme^{-\pi\eta_{ij}/2}\Gamma(1+\rmi\eta_{ij})
  {}_1F_1(-\rmi\eta_{ij};1;\rmi\zeta_{ij})
  \right],
  \label{eq:three-C-wave}\\
  \zeta_{ij}
  &=k_{ij}r_{ij}-\bm k_{ij}\cdot\bm r_{ij},
  \qquad
  \eta_{ij}=\frac{\mu_{ij}g_{ij}}{\hbar^2k_{ij}} .
  \label{eq:three-C-variables}
\end{align}
The pair momenta are constrained by the common three-body kinematics.  The
product reproduces the leading logarithmic phases when all separations are
large, but it is not an exact factorization of the three-body Schr\"odinger
equation.  Cross derivatives and correlations generate subleading terms.

In a cluster sector $\Omega_\alpha$, the close pair must instead be described
by its full two-body state while the spectator feels its net charge and
multipole corrections.  Sectorial expansions by Alt and Mukhamedzhanov and
later refinements construct leading terms in these overlap regimes
\cite{AltMukhamedzhanov:1993,MukhamedzhanovLieber:1996,
Kadyrov:2003ThreeBody,Mukhamedzhanov:2006Leading,Yakovlev:2022Coulomb}.
The literature therefore does contain important three-body Coulomb
asymptotics.  The precise open difficulty is stronger and more useful than
the slogan that \textit{no asymptotic form is known}: there is no analogue of
the single two-body Coulomb wave that is elementary, exact, and uniformly
controlled across all breakup, rearrangement, and coalescence sectors with
simple matching of every scattering amplitude.  The continuing refinement of
terms in different sectors is evidence of this nonuniformity.

The original Faddeev kernel is compact for suitable short-range forces but
loses this property when an unsplit Coulomb potential is inserted.
Faddeev--Merkuriev theory repairs the integral equations by a configuration-
space splitting
\begin{equation}
  V_\alpha^C(\bm x_\alpha)
  =V_\alpha^{(s)}(\bm x_\alpha,\bm y_\alpha)
  +V_\alpha^{(l)}(\bm x_\alpha,\bm y_\alpha),
  \label{eq:merkuriev-splitting}
\end{equation}
chosen so that the short-range part is active in the corresponding two-body
sector and the long-range parts are included in channel Hamiltonians.  The
resulting equations have the compactness and channel structure needed for
analysis and numerical resonance calculations
\cite{FaddeevMerkuriev:1993,Keller:2009FaddeevMerkuriev}.  This is a powerful
solution at the operator-equation level, but it does not reduce the global
coordinate-space boundary condition to a product of three independent
Coulomb waves.

\subsection{From Coulomb phases to soft photons in QED}
\label{subsec:qed-infrared}
% BEGIN CURATED SUBJECT INDEX
\index{superselection sector}
\index{GNS construction}
\index{wave packet}
\index{cross section}
\index{boundary condition!incoming and outgoing}
\index{scattering!Coulomb}
\index{infrared problem}
\index{dressing state}
\index{Fock space!infrared limitation}
\index{observable algebra}
% END CURATED SUBJECT INDEX

The nonrelativistic Coulomb logarithm has a direct field-theoretic analogue.
For massive charged particles moving with asymptotic momenta $p_i$, the
large-time electromagnetic interaction decreases as $1/|t|$, not faster.
Integrating it produces a logarithmic phase and a coherent displacement of
soft photon modes.  The leading dressing function for one charge has the
schematic small-$|\bm k|$ behavior
\begin{equation}
  f_{p,\lambda}(\bm k)
  \sim
  \frac{e\,p\cdot\epsilon_\lambda(\bm k)}
  {(p\cdot k)\sqrt{2|\bm k|}},
  \qquad |\bm k|\to0 ,
  \label{eq:qed-soft-dressing}
\end{equation}
with gauge-dependent longitudinal pieces arranged so that physical
observables satisfy Gauss' law.  Since $f_{p,\lambda}(\bm k)$ behaves as
$|\bm k|^{-3/2}$ for generic direction,
\begin{equation}
  \sum_\lambda\int_{|\bm k|<\Lambda}
  \rmd^3k\,|f_{p,\lambda}(\bm k)|^2
  \sim C\int_0^\Lambda\frac{\rmd|\bm k|}{|\bm k|}
  =\infty .
  \label{eq:soft-photon-number-log}
\end{equation}
The coherent cloud contains an infinite expected number of arbitrarily soft
photons.  Its soft energy can nevertheless remain finite because insertion of
one factor of $|\bm k|$ makes the integral convergent.  The formal Weyl
displacement is therefore not implemented by a unitary operator on the
vacuum Fock space when the infrared cutoff is removed.  This is a change of
representation, not merely a large perturbative coefficient
\cite{BlochNordsieck:1937,Chung:1965,KulishFaddeev:1970}.

The many-particle soft factor depends on the entire scattering process through
the eikonal current
\begin{equation}
  \vsub{J}{soft}^\mu(k)
  =
  \sum_{i\in\mathrm{out}}
  e_i\frac{p_i^\mu}{p_i\cdot k}
  -
  \sum_{j\in\mathrm{in}}
  e_j\frac{p_j^\mu}{p_j\cdot k}.
  \label{eq:eikonal-soft-current}
\end{equation}
Its squared norm contains pairwise terms for every pair of external charges.
Adding particles therefore adds correlated infrared data; it does not merely
tensor another independent one-particle cloud.  Charge conservation removes
some leading pieces, but velocity-dependent soft correlations remain.

Gauss' law sharpens the statement.  A charged state carries an electric flux
detectable at arbitrarily large distances.  Under general assumptions such a
state cannot be a sharp eigenstate of the mass operator
\cite{Buchholz:1986}.  The isolated one-particle delta function expected by
the LSZ/Fock picture is replaced by an infraparticle threshold with continuous
spectral weight.  Thus the phrase \textit{the initial electron state in Fock
space} already fails before an ordinary scattering amplitude is calculated.
Rigorous constructions of infraparticle scattering states exist in models of
nonrelativistic QED and illustrate how much asymptotic information must be
added \cite{ChenFrohlichPizzo:2010}.

Three standard infrared strategies answer different questions:
\begin{enumerate}
  \item Inclusive Bloch--Nordsieck or Yennie--Frautschi--Suura observables sum
  over experimentally unresolved soft radiation.  Virtual and real
  divergences cancel in sufficiently inclusive probabilities
  \cite{YennieFrautschiSuura:1961}.

  \item Chung and Kulish--Faddeev constructions modify the asymptotic states
  and dynamics by coherent photon dressings.  They aim at an infrared-finite
  amplitude between dressed sectors rather than between finite-photon Fock
  vectors \cite{Chung:1965,KulishFaddeev:1970}.

  \item Algebraic approaches retain the observable algebra and describe
  charged or soft sectors by inequivalent representations, asymptotic flux
  observables, and superselection data \cite{Haag:1996,Buchholz:1986}.
\end{enumerate}

The work of Hirai and Sugishita makes the second strategy particularly
concrete and also explains why it cannot be separated from Gauss' law
\cite{HiraiSugishita:2019,HiraiSugishita:2021,HiraiSugishita:2023}.
Write a general asymptotic charged state schematically as
\begin{equation}
  \ket{\alpha;F}
  =\rme^{
    \sum_\lambda\int\rmd^3k\,
    \left(
      F_{\alpha,\lambda}(\bm k)a_\lambda^\dagger(\bm k)
      -F_{\alpha,\lambda}^*(\bm k)a_\lambda(\bm k)
    \right)
  }\vsub{\ket{\alpha}}{hard} .
  \label{eq:general-qed-dressed-state}
\end{equation}
At finite infrared cutoff this is an ordinary coherent state.  Removing the
cutoff is meaningful only after the incoming and outgoing dressings have been
matched to the soft current of the hard process.  In a convention-independent
form, the infrared-safe ``dress code'' is that
\begin{equation}
  \vsup{F_{\beta,\lambda}}{out}(\bm k)
  -\vsup{F_{\alpha,\lambda}}{in}(\bm k)
  -J_{\beta\alpha,\lambda}(\bm k)
  \quad\hbox{is square integrable at $\bm k=0$} ,
  \label{eq:hirai-sugishita-dress-code}
\end{equation}
where $J_{\beta\alpha,\lambda}$ is the transverse polarization component of
the universal eikonal current in Eq.~\eqref{eq:eikonal-soft-current}; signs can
be moved between $J$ and the definitions of the in/out dressings.  Equivalently,
the coefficient of the logarithmically divergent norm must vanish.  Terms
softer than the universal $1/(p\cdot k)$ profile remain genuine choices of the
asymptotic state rather than infrared counterterms.

This condition has three complementary readings.  Dynamically, it is the
cancellation of the virtual soft factor by the coherent clouds attached to
the asymptotic states.  Gauge theoretically, the dressing supplies the
electric field required by the Gauss constraint, so an arbitrary
Faddeev--Kulish-like coherent factor is not automatically sufficient.
Asymptotically, the matching condition is conservation of the large-gauge
charge and can be read as the electromagnetic memory relation between past
and future null infinity.  Hirai and Sugishita further show that the choice of
dressing matters even for inclusive measurements: tracing over unresolved
photons can erase interference between hard wave packets unless the initial
and final clouds obey the corresponding dress code.  The dressed-state and
inclusive prescriptions are therefore compatible but are not interchangeable.

An inclusive cross section can be finite even when no unitary $S$ matrix
exists between the naive charged Fock vectors.  Conversely, specifying a
dressed state requires more information than specifying an experimental
resolution.  Conflating these two statements is a common source of apparent
contradictions in discussions of the infrared problem.

\subsection{Composite charged operators and long-range correlations}
\label{subsec:composite-infrared}
% BEGIN CURATED SUBJECT INDEX
\index{scattering!Coulomb}
\index{scattering!three-body Coulomb}
\index{infrared problem}
\index{dressing state}
% END CURATED SUBJECT INDEX

A local bare electron field is not gauge invariant.  A gauge-invariant charged
operator must also create its electric field.  A Dirac-type dressing is
\begin{equation}
  \Psi_f(x)
  =
  \rme^{
  \rmi e\int\rmd^3y\,
  f^i(\bm x-\bm y)A_i(x^0,\bm y)
  }\psi(x),
  \qquad
  \partial_i f^i(\bm z)=\delta^{(3)}(\bm z).
  \label{eq:dirac-dressed-field}
\end{equation}
The Coulomb choice
$f^i(\bm z)=-\partial^i(4\pi|\bm z|)^{-1}$ extends to spatial infinity and
creates the Coulomb field together with the charge
\cite{Dirac:1955GaugeInvariant}.  A Wilson-line dressing chooses a different
field geometry.  Neither is a compactly supported local composite operator.

It is therefore useful, with a qualification, to describe part of the QED
infrared problem as an infrared singularity of a composite operator produced
by long-range correlations.  The qualification is that three structures must
be distinguished:
\begin{enumerate}
  \item soft divergences of amplitudes between inappropriate Fock states;
  \item the absence of a sharp charged one-particle mass pole, or
  infraparticle phenomenon;
  \item infrared sensitivity of nonlocal gauge-invariant dressings and their
  correlation functions.
\end{enumerate}
They have a common origin in massless propagation and Gauss' law, but they are
not identical divergences.  Wilson-line composites also possess ultraviolet
endpoint or cusp renormalization, which should not be relabeled as infrared
merely because the line reaches infinity.

For several charges, a product of individually dressed fields is generally
not the final asymptotic operator.  Cross terms determined by relative
velocities and total soft charge correlate the dressings.  This mirrors the
three-body Coulomb problem: the long-range cloud is a property of the whole
channel and its cluster limits, not a universal one-body factor attached after
the calculation.

\subsection{Operator-algebraic interpretation of infrared sectors}
\label{subsec:operator-algebra-ir}
% BEGIN CURATED SUBJECT INDEX
\index{measure theory}
\index{direct integral}
\index{von Neumann algebra@\vN\ algebra}
\index{superselection sector}
\index{GNS construction}
\index{wave operator}
\index{boundary condition!incoming and outgoing}
\index{scattering!Coulomb}
\index{Dollard comparison dynamics}
\index{scattering!three-body Coulomb}
% END CURATED SUBJECT INDEX

Let $\vsub{\Alg}{obs}$ be an abstract algebra of gauge-invariant
observables.  The vacuum defines a Fock representation
$\pi_0(\vsub{\Alg}{obs})$ through its GNS state.  A coherent shift by a
square-integrable one-particle wave function is unitarily implemented in that
representation.  The infrared displacement in
Eq.~\eqref{eq:soft-photon-number-log} is not square integrable, so the limiting
charged state can define a representation $\vsub{\pi}{IR}$ disjoint from
$\pi_0$.  The abstract observable algebra can remain the same while the
Hilbert-space representation changes.

Asymptotic electric flux, soft charge, or velocity data may supply central or
superselection labels for an enlarged asymptotic algebra.  Under appropriate
separability and measurability assumptions, a mixed representation can be
organized schematically as
\begin{align}
  \vsub{\Hil}{as}
  &=\int_X^\oplus\Hil_\xi\dd\mu(\xi),
  \label{eq:ir-sector-direct-integral}\\
  \vsub{\vNAlg}{as}
  &=\int_X^\oplus\vNAlg_\xi\dd\mu(\xi).
  \label{eq:ir-algebra-direct-integral}
\end{align}
Here $\xi$ can include soft or flux data and $\mu$ describes the statistical
mixture.  This direct integral is an organizational theorem, not a cure by
itself: a choice of measure and measurable field must be supplied, and
disjoint fibers are not connected by a unitary operator in any one fiber.
Scattering should instead be formulated as an intertwiner between appropriate
incoming and outgoing representations or as a map of asymptotic observable
algebras.

\begin{table}[t]
  \centering
  \small
  \begin{tabular}{>{\raggedright\arraybackslash}p{0.16\textwidth}
                  >{\raggedright\arraybackslash}p{0.23\textwidth}
                  >{\raggedright\arraybackslash}p{0.22\textwidth}
                  >{\raggedright\arraybackslash}p{0.24\textwidth}}
    \toprule
    problem & naive comparison state & long-range obstruction &
    appropriate replacement \\
    \midrule
    short-range two-body &
    free plane wave or packet &
    none after an integrable phase &
    ordinary M{\o}ller wave operators \\
\addlinespace
    two-body Coulomb &
    free plane wave &
    logarithmic phase along a free trajectory &
    Coulomb wave and Dollard dynamics \\
\addlinespace
    three-body Coulomb &
    product of independent pair waves &
    incompatible breakup and cluster limits &
    sectorial asymptotics and Faddeev--Merkuriev channels \\
\addlinespace
    QED charged scattering &
    finite-photon charged Fock vector &
    non-square-integrable soft cloud and Gauss' law &
    inclusive observables, dressed states, or non-Fock sectors \\
    \bottomrule
  \end{tabular}
  \caption{The hierarchy of long-range modifications.  Each row changes the
  asymptotic comparison problem, not the local equations of motion.}
  \label{tab:long-range-hierarchy}
\end{table}

\subsection{Consequences for resonances, complex scaling, and exact WKB}
\label{subsec:long-range-resonances}
% BEGIN CURATED SUBJECT INDEX
\index{spectrum!point}
\index{GNS construction}
\index{Jost function}
\index{boundary condition!incoming and outgoing}
\index{resonance!residue}
\index{scattering!Coulomb}
\index{infrared problem}
\index{dressing state}
\index{exact WKB}
\index{complex scaling}
% END CURATED SUBJECT INDEX

Long-range asymptotics changes the analytic surface on which a resonance is
defined.  Because $\eta\propto1/k$, the Coulomb factors contain
$\Gamma(l+1+\rmi\eta)$ and logarithms with an essential threshold structure
at $k=0$.  A short-range square-root sheet is therefore insufficient.  Pole
searches should factor the pure Coulomb contribution and continue the
Coulomb-modified Jost determinant.  Response functions should be compared
with a reference Hamiltonian carrying the same Coulomb tail.

Complex scaling can be applied to important Coulombic Hamiltonians under
appropriate dilation-analytic hypotheses, but it does not restore free
asymptotics.  Exterior scaling, channel-adapted coordinates, and
Faddeev--Merkuriev decompositions may be required in many-body problems.  A
rotated eigenvalue must be interpreted relative to the correct Coulomb
thresholds and reference continuum.

Exact WKB displays the two-body logarithm immediately.  At large $r$,
\begin{equation}
  p(r)
  =\sqrt{2\mu\left(E-\frac{g}{r}\right)}
  =\hbar k-\frac{\mu g}{\hbar k}\frac{1}{r}
  +O(r^{-2}),
  \label{eq:coulomb-WKB-momentum}
\end{equation}
so
\begin{equation}
  \frac{1}{\hbar}\int^r p(r')\dd r'
  =kr-\eta\ln r+O(r^{-1}).
  \label{eq:coulomb-WKB-log}
\end{equation}
The Coulomb phase is therefore a global WKB period reaching infinity.  In the
three-body problem the corresponding Hamilton--Jacobi function is
multivariable and its cluster limits are nonuniform.  In QED the asymptotic
phase becomes an operator-valued soft dressing.

A genuinely field-theoretic resonance theory must combine these infrared
representations with analytic continuation of multiparticle correlation
functions.  For a charged unstable excitation, an infraparticle cut and a
resonance pole need not be separable by the elementary pole-plus-background
ansatz.  Constructing an asymptotic algebra, its incoming and outgoing GNS
representations, and a continued spectral object carrying widths and residues
is a central open problem in extending the present program from potential
scattering to gauge theory.
